\chapter{Řízení servomotoru pomocí NE555}
\label{chap:ne555}

\section{Princip řízení}

Pro ruční ovládání servomotoru použijeme časovač NE555P napájený 5 V. V astabilním zapojení vytváří pravidelné pulzy, jejichž délku nastavujeme potenciometrem. K řízení tak není potřeba mikrokontrolér ani program.~\cite{servo555prakticky,ne555}

Princip spočívá v opakovaném nabíjení a vybíjení kondenzátoru \textbf{C1} (viz schéma \ref{fig:ne555}) mezi přibližně třetinou a dvěma třetinami napájecího napětí. Při nabíjení přes \textbf{R1}, \textbf{P1} a diodu \textbf{D1} je výstup v logické jedničce. Po dosažení horního prahu se přepne do nuly a kondenzátor se vybíjí přes \textbf{R2} a vývod 7. Dioda odděluje obě cesty, čímž umožňuje vytvořit krátký kladný pulz a~dlouhou mezeru.~\cite{ne555}

\section{Výpočet časování a volba součástek}

Požadujeme pulzy o délce 0,5--2,5 ms s periodou přibližně 20 ms. Délku kladného pulzu $t_H$ určují \textbf{R1}, potenciometr \textbf{P1} a \textbf{C1}, zatímco délku mezery $t_L$ určuje především \textbf{R2} a \textbf{C1}. Při uvažovaném úbytku napětí na diodě $U_D$ platí
\begin{align}
 t_H&\approx(R_1+R_P)C_1
 \ln\!\left(\frac{2U_{CC}/3-U_D}{U_{CC}/3-U_D}\right),\\
 t_L&\approx\ln(2)R_2C_1,\\
 T&=t_H+t_L,\qquad f=\frac{1}{T}.
\end{align}

Volíme \textbf{R1} = 3,3~k\si{\ohm}, \textbf{P1} = 25~k\si{\ohm}, \textbf{R2} = 221~k\si{\ohm}, \textbf{C1} = 120 nF a \textbf{D1} = 1N4148. Při napájení 5~V a úbytku na diodě 0,6 V vychází rozsah pulzů přibližně 0,37--3,20 ms a délka mezery 18,38~ms.

Při omezení pulzů na požadovaných 0,5--2,5 ms se perioda mění mezi 18,88 a 20,88 ms, tedy přibližně 48--53 Hz. Perioda proto není zcela konstantní, což ale v případě serva nevadí. Skutečné časování ovlivňují také tolerance součástek, a proto pokud je potřeba přesné řízení, je vhodné provést doladění s osciloskopem.

\section{Schéma zapojení}

Navržené zapojení časovače NE555 je znázorněno na obrázku~\ref{fig:ne555}.

\noindent\begin{minipage}{\textwidth}\centering
\begin{tikzpicture}[transform shape]
	% Paths, nodes and wires:
	\draw[line width=1.2pt,fill=light] (7, 2.5) -- (11.25, 2.5) -| (11.25, 7) -- (7, 7) -| (7, 2.5);
	\node[shape=rectangle, minimum width=0.965cm, minimum height=0.365cm] at (9.75, 5.05){} node[anchor=center, align=center, text width=0.929cm, inner sep=1pt] at (9.75, 5.05){\small U1};
	\node[shape=rectangle, minimum width=1.765cm, minimum height=0.465cm] at (9.85, 4.5){} node[anchor=center, align=center, text width=1.729cm, inner sep=1pt] at (9.85, 4.5){\small NE555P};
	\draw (7, 6.3) -- (3, 6.3);
	\draw (7, 4.5) -- (6.2, 4.5) -| (6.2, 3.4);
	\draw (7, 3.4) -- (3, 3.4);
	\node[shape=rectangle, minimum width=1.465cm, minimum height=0.315cm] at (7.9, 6.3){} node[anchor=west, align=left, text width=1.429cm, inner sep=1pt] at (7.15, 6.3){\scriptsize 7 DISCH};
	\node[shape=rectangle, minimum width=1.465cm, minimum height=0.315cm] at (7.9, 4.5){} node[anchor=west, align=left, text width=1.429cm, inner sep=1pt] at (7.15, 4.5){\scriptsize 6 THRES};
	\node[shape=rectangle, minimum width=1.465cm, minimum height=0.315cm] at (7.9, 3.4){} node[anchor=west, align=left, text width=1.429cm, inner sep=1pt] at (7.15, 3.4){\scriptsize 2 TRIG};
	\draw (11.25, 4.75) |- (12.3, 4.8);
	\node[shape=rectangle, minimum width=2.465cm, minimum height=0.815cm] at (13.7, 4.8){} node[anchor=west, align=left, text width=2.429cm, inner sep=1pt] at (12.45, 4.8){\small 3 OUT\\signál pro servo};
	\draw (8, 7) |- (7.75, 11) -| (9.5, 7.5) -| (9.5, 7);
	\node[shape=rectangle, minimum width=0.965cm, minimum height=0.265cm] at (8, 6.75){} node[anchor=center, align=center, text width=0.929cm, inner sep=1pt] at (8, 6.75){\scriptsize 8 VCC};
	\node[shape=rectangle, minimum width=1.865cm, minimum height=0.115cm] at (9.8, 6.675){} node[anchor=center, align=center, text width=1.829cm, inner sep=1pt] at (9.8, 6.675){\scriptsize 4 RESET};
	\draw (8, 2.5) -| (8, 1.5);
	\node[ground] at (8, 0.5){};
	\node[shape=rectangle, minimum width=1.165cm, minimum height=0.315cm] at (8, 2.75){} node[anchor=center, align=center, text width=1.129cm, inner sep=1pt] at (8, 2.75){\scriptsize 1 GND};
	\draw (9.5, 2.5) to[capacitor] (9.5, 0.6);
	\node[ground] at (9.5, 0.5){};
	\node[shape=rectangle, minimum width=1.265cm, minimum height=0.315cm] at (9.5, 2.75){} node[anchor=center, align=center, text width=1.229cm, inner sep=1pt] at (9.5, 2.75){\scriptsize 5 CONT};
	\node[shape=rectangle, minimum width=1.965cm, minimum height=0.465cm] at (11, 1.55){} node[anchor=center, align=center, text width=1.929cm, inner sep=1pt] at (11, 1.55){\small C7 10 nF};
	\node[shape=rectangle, minimum width=1.465cm, minimum height=0.465cm] at (3, 11.789){} node[anchor=center, align=center, text width=1.429cm, inner sep=1pt] at (3, 11.789){\small +5 V};
	\draw (3, 11) to[european resistor] (3, 9.2);
	\draw (3, 9.2) -| (3, 8.8);
	\node[shape=rectangle, minimum width=2.465cm, minimum height=0.465cm] at (1.35, 10.1){} node[anchor=center, align=center, text width=2.429cm, inner sep=1pt] at (1.35, 10.1){\small R1 3,3 kΩ};
	\draw (2.75, 8.8) -- (3.25, 8.8) -| (3.25, 7.3) -- (2.75, 7.3) -| (2.75, 8.8);
	\node[shape=rectangle, minimum width=2.265cm, minimum height=0.815cm] at (1.3, 8.05){} node[anchor=center, align=center, text width=2.229cm, inner sep=1pt] at (1.3, 8.05){\small P1\\25 kΩ lin.};
	\node[shape=rectangle, minimum width=0.365cm, minimum height=0.265cm] at (2.65, 9){} node[anchor=center, align=center, text width=0.329cm, inner sep=1pt] at (2.65, 9){\scriptsize A};
	\node[shape=rectangle, minimum width=0.365cm, minimum height=0.265cm] at (2.65, 7.05){} node[anchor=center, align=center, text width=0.329cm, inner sep=1pt] at (2.65, 7.05){\scriptsize B};
	\draw[-latex] (4.5, 8.05) -- (3.25, 8.05);
	\node[shape=rectangle, minimum width=0.415cm, minimum height=0.315cm] at (4.85, 8.05){} node[anchor=center, align=center, text width=0.379cm, inner sep=1pt] at (4.85, 8.05){\scriptsize W};
	\draw (4.5, 8.05) -| (4.5, 7) -- (3, 7);
	\draw (3, 7.3) -| (3, 6.3);
	\draw (3, 6.3) to[european resistor] (3, 3.4);
	\node[shape=rectangle, minimum width=2.565cm, minimum height=0.465cm] at (1.3, 4.85){} node[anchor=center, align=center, text width=2.529cm, inner sep=1pt] at (1.3, 4.85){\small R2 221 kΩ};
	\draw (5.1, 6.3) to[empty diode] (5.1, 3.4);
	\node[shape=rectangle, minimum width=1.165cm, minimum height=0.615cm] at (5.9, 5.551){} node[anchor=center, align=center, text width=1.129cm, inner sep=1pt] at (5.9, 5.551){\scriptsize D1\\1N4148};
	\node[shape=rectangle, minimum width=0.365cm, minimum height=0.265cm] at (4.7, 5.8){} node[anchor=center, align=center, text width=0.329cm, inner sep=1pt] at (4.7, 5.8){\scriptsize A};
	\node[shape=rectangle, minimum width=0.365cm, minimum height=0.265cm] at (4.7, 4){} node[anchor=center, align=center, text width=0.329cm, inner sep=1pt] at (4.7, 4){\scriptsize K};
	\draw (3, 3.4) to[capacitor] (3, 1.2);
	\node[ground] at (3, 0.5){};
	\node[shape=rectangle, minimum width=2.565cm, minimum height=0.465cm] at (1.3, 2.3){} node[anchor=center, align=center, text width=2.529cm, inner sep=1pt] at (1.3, 2.3){\small C1 120 nF};
	\node[circ] at (3, 7){};
	\node[circ] at (3, 6.3){};
	\node[circ] at (5.1, 6.3){};
	\node[circ] at (3, 3.4){};
	\node[circ] at (5.1, 3.4){};
	\node[circ] at (6.2, 3.4){};
	\draw (10.75, 10.95) to[capacitor] (10.75, 8.75);
	\node[ground] at (10.75, 8.75){};
	\draw (13.45, 10.95) to[capacitor] (13.45, 8.75);
	\node[ground] at (13.45, 8.75){};
	\node[shape=rectangle, minimum width=2.165cm, minimum height=0.465cm] at (11.65, 9.25){} node[anchor=center, align=center, text width=2.129cm, inner sep=1pt] at (11.65, 9.25){\small C5 100 nF};
	\node[shape=rectangle, minimum width=2.165cm, minimum height=0.465cm] at (14.4, 9.25){} node[anchor=center, align=center, text width=2.129cm, inner sep=1pt] at (14.4, 9.25){\small C6 10 µF};
	\node[shape=rectangle, minimum width=0.365cm, minimum height=0.365cm] at (13, 10.2){} node[anchor=center, align=center, text width=0.329cm, inner sep=1pt] at (13, 10.2){\small +};
	\node[shape=rectangle, minimum width=4.365cm, minimum height=0.765cm] at (12.056, 7.519){} node[anchor=center, align=center, text width=4.329cm, inner sep=1pt] at (12.056, 7.519){\scriptsize C5 umístit co nejblíže\\napájecím vývodům NE555};
	\draw (8, 1.5) -| (8, 0.5);
	\draw (9.5, 0.6) -| (9.5, 0.5);
	\draw (3, 1.2) -- (3, 0.5);
	\draw (3, 11) -| (13.45, 10.95);
	\draw (10.75, 10.95) -| (10.75, 11);
	\node[vcc] at (3, 11){};
	\node[circ] at (8, 11){};
	\node[circ] at (9.5, 11){};
	\node[circ] at (10.75, 11){};
	\node[circ] at (3, 11){};
\end{tikzpicture}
\captionsetup{hypcap=false}\captionof{figure}{Generátor řídicích pulzů s~NE555.}
\label{fig:ne555}
\end{minipage}\medskip

Vývody 2 (TRIG) a 6 (THRES) sledují napětí na \textbf{C1}. Dioda \textbf{D1} má anodu na vývodu 7~(DISCH) a katodu u \textbf{C1}. Jezdec potenciometru \textbf{P1} je spojen s jeho koncem B, čímž se omezuje riziko úplného přerušení odporové dráhy při ztrátě kontaktu.

Vývod 4 (RESET) je připojen na napájení +5 V, vývod 5 (CONT) filtruje \textbf{C7} = 10 nF. Napájení blokují \textbf{C5} = 100 nF a \textbf{C6} = 10~\si{\micro\farad}. Výstup 3 (OUT) vede přes \textbf{R6} = 330~\si{\ohm} na signál serva, zatímco \textbf{R7} = 100~k\si{\ohm} zajišťuje nízkou úroveň při odpojeném generátoru. ~\cite{ne555}